% Working Discussion, transferred without substantive revision.
\section{Discussion}
\label{sec:discussion}

The central purpose of this study is to identify a practical procedure for constructing a large
conformational dataset with coverage of experimentally observed ligand geometries. The benchmark
separates several properties that are often combined under the general description of conformer
quality. Structural validity determines whether a sample satisfies the specified identity and
geometry requirements. Geometric diversity describes the variation represented in an ensemble.
Energy analysis characterizes that variation under a common molecular model, while
experimental-conformation recovery assesses whether the ensemble includes shapes observed in protein
complexes.

The preliminary results support the value of considering these properties together. Broad geometric
exploration does not guarantee recovery, and strong recovery by the closest member of an ensemble
does not establish that most samples lie near observed conformations. The comparison between
single-reference recovery and multi-reference coverage makes this distinction explicit. It also
explains why the preferred generator may depend on whether the immediate objective is to recover at
least one useful conformation, cover several experimental conformations, or concentrate sampling
near those observations.

These evaluations inform dataset construction without establishing that every generated conformer is
biologically relevant. The experimental collections sample only part of the accessible
conformational space, and the force-field calculations do not represent binding free energies. The
results are also conditional on the selected molecules, generation settings, and sampling budgets.
Conclusions about Qwen's generalization will depend on the evaluation after benchmark molecules have
been excluded from its training data.

The proposed public release is intended to provide a reproducible resource for learning molecular
geometry, particularly for pharmacophore-conditioned generation. Its demonstrated properties will be
the structural coverage, diversity, and validity measured in this study. Whether training on this
resource improves a particular downstream model remains a separate question for future work.
